\documentclass[12pt]{article}
\usepackage[utf8]{inputenc}
\usepackage[T1]{fontenc}
\usepackage[margin=1in]{geometry}
\usepackage{setspace}
\usepackage{xcolor}
\usepackage[colorlinks=true, urlcolor=blue!60!black, linkcolor=blue!60!black]{hyperref}
\setlength{\parindent}{0em}
\setlength{\parskip}{0.6em}
\renewcommand{\baselinestretch}{1.0}
\begin{document}
\begin{center}
{\LARGE \textit{Miller v. California}, 413 U.S. 15 (1973)\par}
\vspace{0.5em}
{\large Case Brief --- Two-Page Edition\par}
\end{center}

\section{Memory Jogger}
Obscenity remains unprotected, but the Court replaced the ``utterly without redeeming social value'' test with a three-prong test keyed to local community standards and the absence of serious literary, artistic, political, or scientific value (\href{https://supreme.justia.com/cases/federal/us/413/15/}{413 U.S. 15}).

\section{Detailed Case Facts}
Marvin Miller's mail-order business sent five unsolicited brochures, consisting mainly of explicit pictures of group sexual activity, to a Newport Beach restaurant, whose manager complained to police. He was convicted under California Penal Code \S~311.2(a), which defined obscenity along the lines of the \textit{Memoirs v. Massachusetts} plurality.

\section{Procedural History}
The jury was instructed on statewide community standards; the Appellate Department of the Superior Court summarily affirmed, and the Supreme Court vacated and remanded.

\section{Judicial Votes}
Vacated, 5--4. Burger wrote for White, Blackmun, Powell, and Rehnquist; Douglas dissented alone; Brennan, joined by Stewart and Marshall, dissented.

\section{Holding}
Obscenity is unprotected. A state may regulate works depicting sexual conduct specifically defined by state law if (a) the average person, applying contemporary community standards, finds the work as a whole appeals to the prurient interest; (b) it depicts that conduct in a patently offensive way; and (c) taken as a whole, it lacks serious literary, artistic, political, or scientific value. No national standard is required.

\section{Analysis of Opinions}
Burger sought to end the Court's role as ``an unreviewable board of censorship for the 50 States.'' The \textit{Memoirs} plurality forced prosecutors to prove a negative, and a national standard would be ``an exercise in futility'' in a diverse nation. The new test isolates hard-core material and gives fair notice through specific statutory definitions. Douglas argued that ``obscenity'' is absent from the Constitution, that retroactive standards carry ``all the evils of an \textit{ex post facto} law,'' and that an ``offensive'' standard guts the First Amendment. Brennan concluded that no formulation could separate obscene from protected speech and found the statute overbroad.

\section{Examples: Future Applications}
\textbf{Same-Side:} unsolicited pamphlets of actual intercourse with no commentary, mailed into a county whose statute enumerates those acts, are unprotected. \textbf{Opposite-Side:} a medical textbook with graphic anatomical photographs has serious scientific value and is protected. \textbf{Fence-Sitter:} an explicit online film with a political monologue prosecuted in a rural district.

\section{Critique}
Defenders cite historical restrictions and praise the federalism of local standards; textualists echoing Douglas see an unmoored invention. Civil libertarians fault its vagueness and chilling effect, magnified online, while Catharine MacKinnon argued it targets moral offense rather than pornography's harm to women.

\section{Key Quotations}
\begin{quote}
``The basic guidelines for the trier of fact must be: (a) whether `the average person, applying contemporary community standards' would find that the work, taken as a whole, appeals to the prurient interest, \textit{Kois v. Wisconsin}, supra, at 230, quoting \textit{Roth v. United States}, supra, at 489; (b) whether the work depicts or describes, in a patently offensive way, sexual conduct specifically defined by the applicable state law; and (c) whether the work, taken as a whole, lacks serious literary, artistic, political, or scientific value.''
\end{quote}
\noindent --- Chief Justice Burger, Opinion of the Court, 413 U.S. at 24.

\end{document}
